\documentclass[10pt,a4paper]{article}
\usepackage{amsmath,amssymb,amsthm,graphicx}
\usepackage{amsmath,amssymb,amsthm,hyperref,xfrac}
\providecommand{\ie}{i.e.\ }
\providecommand{\skpt}{\smallskip}
\usepackage[scr=rsfs,cal=euler]{mathalfa}
\multlinegap0pt

\hyphenation{begin decom-po-si-tion sequence}
\newenvironment{enumeratea}
{\bgroup\def\theenumi{\alph{enumi}}\begin{enumerate}}
{\end{enumerate}\egroup}

\newcommand{\ba}{{\boldsymbol a}}
\newcommand{\bb}{{\boldsymbol b}}
\newcommand{\bd}{{\boldsymbol d}}
\newcommand{\bu}{{\boldsymbol u}}
\newcommand{\bx}{{\boldsymbol x}}
\newcommand{\bmy}{{\boldsymbol y}}

\newcommand{\bD}{{\boldsymbol D}}
\newcommand{\bL}{{\boldsymbol L}}
\newcommand{\bP}{{\boldsymbol P}}
\newcommand{\bV}{{\boldsymbol V}}
\newcommand{\bX}{{\boldsymbol X}}
\newcommand{\bZ}{{\boldsymbol Z}}

\newcommand{\Psfrac}[2]{(\sfrac{#1}{#2})}
\newcommand{\psfrac}[2]{\sfrac{(#1)}{#2}}
\newcommand{\bigO}[1]{O(#1)}

\usepackage{mathtools}
\usepackage{bbm}

\newcommand{\NN}{\mathbb{N}}
\newcommand{\ZZ}{\mathbb{Z}}
\newcommand{\RR}{\mathbb{R}}
\newcommand{\CC}{\mathbb{C}}
\newcommand{\PP}{\mathbb{P}}
\newcommand{\EE}{\mathbb{E}}

\newcommand{\boldalpha}{\boldsymbol{\alpha}}
\newcommand{\bolddelta}{\boldsymbol{\delta}}
\newcommand{\boldrho}{\boldsymbol{\rho}}
\newcommand{\dee}{\,\mathrm{d}}
\newcommand{\one}{\mathbbm{1}}
\newcommand{\BC}{\mathcal{B}}
\newcommand{\EC}{\mathcal{E}}
\newcommand{\RS}{\mathscr{R}}
\newcommand{\prob}[1]{\PP[ #1 ]}
\newcommand{\bigoo}[1]{O\!\left(#1 \right)}
\newcommand{\abs}[1]{\left\lvert #1 \right\rvert}
\newcommand{\floor}[1]{\left\lfloor #1 \right\rfloor}
\newcommand{\riemint}[4]{\int_{#1}^{#2}\! #3 \, \mathrm{d}#4}
\newcommand{\maxnorm}[1]{\left\| #1 \right\|_{\infty}}

\renewcommand{\Re}{\textup{Re}}
\renewcommand{\Im}{\textup{Im}}

\newcommand{\dtv}{\mathrm{d}_{\mathrm{TV}}}
\newcommand{\dw}{\mathrm{d}_{\mathrm W}}

\let\epsilon\varepsilon

\newcommand{\pextra}{P_{\mathtt{extra}}}
\newcommand{\parratia}{P_{\mathtt{Arratia}}}

\newcommand{\simplex}{\Delta^{k-1}}

\newcommand{\bg}{\big}
\newcommand{\bgg}{\Big}
\newcommand{\bggg}{\bigg}

\newcommand{\Rnt}{R_{\mathrm{NT}}}
\newcommand{\Rprob}{R_{\mathrm{PD}}}
\newcommand{\ECnt}{\EC_{\mathrm{NT}}}
\newcommand{\ECprob}{\EC_{\mathrm{PD}}}

\newcommand{\Bnt}{\mathcal{B}_{\mathrm{NT}}}
\newcommand{\Bprob}{\mathcal{B}_{\mathrm{PD}}}

\DeclareMathOperator{\vard}{Var}
\DeclareMathOperator{\Dir}{Dir}

\theoremstyle{plain}
\newtheorem{cor}{Corollary}[section]
\newtheorem{lem}[cor]{Lemma}
\newtheorem{prop}[cor]{Proposition}
\newtheorem{conj}[cor]{Conjecture}
\newtheorem{theo}{Theorem}

\theoremstyle{definition}
\newtheorem{example}{Example}
\newtheorem{deff}[cor]{Definition}
\newtheorem{rem}{Remark}[section]
\newtheorem*{rems}{Remark}
\newtheorem*{remss}{Remarks}

\makeatletter
\ams@newcommand{\multiint}[1]{\DOTSI\protect\MultiIntegral{#1}}
\renewcommand{\MultiIntegral}[1]{%
\edef\ints@c{\noexpand\intop
\ifnum#1=\z@\noexpand\intdots@\else\noexpand\intkern@\fi
\replicate{#1-2}{\noexpand\intop\noexpand\intkern@}%
\noexpand\intop
\noexpand\ilimits@
}%
\futurelet\@let@token\ints@a
}
\makeatother
\ExplSyntaxOn
\cs_new:Npn \replicate #1 #2 { \prg_replicate:nn { #1 } { #2 } }
\ExplSyntaxOff
\title{When the three good events hold, the integer-friendly random integer coincides with Arratia's: J\_x = J\_x\textasciicircum{}*}
\author{Tony Haddad, Dimitris Koukoulopoulos}
\date{doi:10.5802/jep.317; CC BY 4.0}
\begin{document}
\maketitle
\noindent\textit{Editorial note.} The counted statement is the author's labelled lemma `lem:decomposition J neq J*` (source lines 842-848); the proof environment begins on the next line after \textbackslash{}end\{lem\} apart from one blank line, so the pairing is by adjacency.\par\medskip
\noindent\textit{Editorial note.} The retained context is the opening of the author's own section on the integer-friendly version of J\_x, which defines the von Mangoldt weight h that both the statement and the proof use.\par\medskip
\noindent\textit{Editorial note.} No counter adjustment is retained: the publisher numbering is not reproduced and the item prints with the wrapper's own sequential numbering.\par\medskip
\noindent\textit{Editorial note.} No retained line refers to a label, citation or figure outside the retained span, so the package declares no external reference.\par\medskip
\noindent\textit{Editorial scope.} One original result of the pinned publisher article, counted once: the event decomposition forcing J\_x = J\_x\textasciicircum{}* (author label lem:decomposition J neq J*).\par\medskip
\noindent\textit{Allowed standard prerequisites (editorial).}\begin{itemize}
\item The two random integers. J\_x = prod\_\{i>=1\} e\textasciicircum{}\{h(Y\_i)\} is the integer-friendly product, and J\_x\textasciicircum{}* is Arratia's random integer defined by the ordered indexing of the transformed point process R\textasciicircum{}*.. The statement compares the two and the proof computes both products over the index sets B\_+ and B\_-. 
\item The good events E\_1, E\_2, E\_3. E\_1, E\_2, E\_3 are the events displayed immediately before the lemma, controlling the spacing sums S\_1, S\_0, the error R\_x and the ordering of the ratios W\_i Y\_i / h(Y\_i).. Each event is consumed at a different place: E\_1 at lines 850-864, E\_2 and E\_3 at lines 866-877.
\item von Mangoldt function. Lambda(n) = log p if n is a prime power p\textasciicircum{}k and 0 otherwise.. It is the weight h used in the products J\_x and in the section's construction, retained in the context.
\end{itemize}\medskip
\section{An integer-friendly version of \texorpdfstring{$J_x$}{Jx}}
\label{sec:integer friendly J}

Let $\Lambda$ be the von Mangoldt function, that is to say
\[
\Lambda(n) = \begin{cases}
\log p &\text{if $n = p^k$ for some prime power $p^k$,} \\
0 &\text{otherwise.}
\end{cases}
\]

\begin{lem}
\label{lem:decomposition J neq J*}
For $x > 1$, we~have $\EC_1\cap \EC_2\cap \EC_3\subseteq \{ J_x = J_x^*\}$. In~particular,
\[
\PP\bg[J_x \ne J_x^*\bg] \le \PP[\EC_1^c] + \PP[\EC_1 \cap \EC_2^c] + \PP[\EC_1 \cap \EC_3^c].
\]
\end{lem}

\begin{proof} Recall that $J_x=\prod_{i\ge1}e^{h(Y_i)}$ and assume that $\mathcal E_1$ occurs. Then
\begin{equation}
\label{eq:e1 implies J < x}
\log J_x = \sum_{i=1}^{\infty} h(Y_i) \le \sum_{i \ge 1} \bg(Y_i + r(Y_i)\bg) \le S_1 + R_x +r_0 < \log x,
\end{equation}
since $r(Y_1)\le r_0$. Similarly, we~have
\[
\log\!\bg(J_x\cdot e^{h(Y_0)}\bg) = \sum_{i=0}^{\infty} h(Y_i) \ge \sum_{i \ge 0} \bg(Y_i - r(Y_i)\bg) \ge S_0 - R_x -2r_0 > \log x.
\]
Therefore, $\EC_1$ implies the inequalities
\begin{equation}
\label{eq:J_x inequalities under E_1}
J_x < x < J_x e^{h(Y_0)}.
\end{equation}
In particular, we~must have $Y_0>e^{-\gamma}$.

Assume now further that $\EC_2$ and $\EC_3$ also occur. We~claim that this implies $J_x = J_x^*$. Let $B_+$ be the set of integers $i \ge 1$ such that $Y_i > e^{-\gamma}$, and let $B_-$ be the set of integers $i\le-1$ such that $Y_i>e^{-\gamma}$.
By our assumption that $\EC_2\cap\EC_3$ occurs and by the definition of $\eta$, we~have
\begin{equation}
\label{eq:preserving order 1}
\frac{W_iY_i}{h(Y_i)} > \frac{W_0Y_0}{h(Y_0)} \qquad\text{if}\ i\in B_+
\end{equation}
and, similarly,
\begin{equation}
\label{eq:preserving order 2}
\frac{W_iY_i}{h(Y_i)} < \frac{W_0Y_0}{h(Y_0)} \qquad\text{if}\ i\in B_-.
\end{equation}
Now, let $K \ge 0$ and $(T_j^*,Q_j^*)$ with $j\in\ZZ_{\le K}$ be the indexing of points of $\RS^*$ given in Section \ref{sec:integer friendly J}, that is to say we have $T_{j-1}^*<T_{j}^*$ for all $j \in \ZZ_{\le K}$, and such that $\prod_{j=0}^K Q_j^*>x \ge \prod_{j=1}^K Q_j^*$. Let $j_0\in\ZZ_{\le K}$ be such that $T_{j_0}^* = W_0Y_0/h(Y_0)$ (which exists because $Y_0>e^{-\gamma}$). Using relations \eqref{eq:preserving order 1} and \eqref{eq:preserving order 2}, we~find
\[
\prod_{j_0 < j \le K} Q_j^*=\prod_{i\in B_+} e^{h(Y_i)} = \prod_{i\ge1} e^{h(Y_i)} = J_x.
\]
On the other hand, we~have $\prod_{j_0 \le j \le K} Q_j^*=e^{h(Y_0)} J_x$. Hence, using \eqref{eq:J_x inequalities under E_1}, we~find that $\prod_{j_0 \!<\! j \!\le\! K}Q_j^*\!<\!x\!<\!\prod_{j_0 \le j \le K}Q_j^*$. In~particular, $j_0\!=\!0$ and thus $J_x^*=\prod_{j_0 < j \le K} Q_j^*\!=\!J_x$, as claimed.
\end{proof}

\begin{thebibliography}{99}
\bibitem{Arratia98} R. Arratia - “On the central role of scale invariant Poisson processes on (0,∞)”, in Microsurveys in discrete probability (Princeton, NJ, 1997), DIMACS Ser. Discrete Math. Theoret. Comput. Sci., vol. 41, American Mathematical Society, Providence, RI, 1998, p. 21-41
\bibitem{Arratia02} R. Arratia - “On the amount of dependence in the prime factorization of a uniform random integer”, in Contemporary combinatorics, Bolyai Soc. Math. Stud., vol. 10, János Bolyai Math. Soc., Budapest, 2002, p. 29-91
\bibitem{ABT06} R. Arratia, A. D. Barbour \& S. Tavaré - “A tale of three couplings: Poisson-Dirichlet and GEM approximations for random permutations”, Combin. Probab. Comput. 15 (2006) no. 1-2, p. 31-62
\bibitem{BareikisManstavicius07} G. Bareikis \& E. Manstavičius - “On the DDT theorem”, Acta Arith. 126 (2007) no. 2, p. 155-168
\bibitem{Bill72} P. Billingsley - “On the distribution of large prime divisors”, Period. Math. Hungar. 2 (1972), p. 283-289
\bibitem{delaBretecheTenenbaum16} R. de la Bretèche \& G. Tenenbaum - “Sur les processus arithmétiques liés aux diviseurs”, Adv. in Appl. Probab. 48 (2016) no. A, p. 63-76
\bibitem{DDT79} J.-M. Deshouillers, F. Dress \& G. Tenenbaum - “Lois de répartition des diviseurs. I”, Acta Arith. 34 (1979) no. 4, p. 273-285 (loose errata)
\bibitem{DonnellyTavare87} P. Donnelly \& S. Tavaré - “The population genealogy of the infinitely-many neutral alleles model”, J. Math. Biol. 25 (1987) no. 4, p. 381-391
\bibitem{Feng10} S. Feng - The Poisson-Dirichlet distribution and related topics. Models and asymptotic behaviors, Probability and its Appl., Springer, Heidelberg, 2010
\bibitem{JKB} N. L. Johnson, S. Kotz \& N. Balakrishnan - Discrete multivariate distributions, Wiley Series in Probab. and Stat.: Applied Probab. and Stat., John Wiley \& Sons, Inc., New York, 1997
\bibitem{Kingman93} J. F. C. Kingman - Poisson processes, Oxford Studies in Probab., vol. 3, The Clarendon Press, Oxford University Press, New York, 1993
\bibitem{KotzBal00} S. Kotz, N. Balakrishnan \& N. L. Johnson - Continuous multivariate distributions. Vol. 1. Models and applications, Wiley Series in Probab. and Stat.: Applied Probab. and Stat., Wiley-Interscience, New York, 2000
\bibitem{Koukoulo19} D. Koukoulopoulos - The distribution of prime numbers, Graduate Studies in Math., vol. 203, American Mathematical Society, Providence, RI, 2019
\bibitem{LastPenrose18} G. Last \& M. Penrose - Lectures on the Poisson process, Institute of Math. Stat. Textbooks, vol. 7, Cambridge University Press, Cambridge, 2018
\bibitem{Ken22} S.-K. Leung - “Dirichlet law for factorisation of integers, polynomials and permutations”, Math. Proc. Cambridge Philos. Soc. 175 (2023) no. 3, p. 649-676
\bibitem{KenWeb24} S.-K. Leung - “Dirichlet law with arbitrary parameters for factorizations of integers” (2025), Retrieved on 24 August 2025 from https://sites.google.com/view/sunkaileung/miscellaneous
\bibitem{MV07} H. L. Montgomery \& R. C. Vaughan - Multiplicative number theory. I. Classical theory, Cambridge Studies in Advanced Math., vol. 97, Cambridge University Press, Cambridge, 2007
\bibitem{NyandwiSmati13} S. Nyandwi \& A. Smati - “Distribution laws of pairs of divisors”, Integers 13 (2013), article ID A13, 13 pages
\bibitem{Tenenbaum00} G. Tenenbaum - “A rate estimate in Billingsley’s theorem for the size distribution of large prime factors”, Q. J. Math. 51 (2000) no. 3, p. 385-403
\end{thebibliography}
\end{document}
